\subsection{The Open Mapping Theorem}\label{sec:open-mapping-theorem}

Let $X$ and $Y$ be Banach spaces. The following three theorems on operators from $X$ to $Y$ have important consequences in our attempts to invert operators.

\begin{Theorem}{Open mapping theorem}
\label{thm:open-mapping}
Let $A\in\mathcal{L}(X,Y)$, and suppose $A$ is surjective (onto). Then the image of every open set $S\subset X$ is open in $Y$.
\end{Theorem}

\begin{Theorem}{Bounded inverse theorem}
\label{thm:bounded-inverse}
Let $A\in\mathcal{L}(X,Y)$, and suppose $A$ is bijective (one-to-one and onto). Then the inverse map $A^{-1}$ is bounded.
\end{Theorem}

\begin{Theorem}{Closed graph theorem}
\label{thm:closed-graph}
Let $(\mathcal{D}(A),A)$ be a linear operator from $X$ to $Y$. Then if $(\mathcal{D}(A),A)$ is a closed operator and its domain $\mathcal{D}(A)$ is closed in $X$, then the operator is bounded.
\end{Theorem}

In a typical course in functional analysis, the open mapping theorem is proved using the Baire category theorem, and the bounded inverse theorem and closed graph theorem are derived as consequences. In fact, as we show in Lemma~\ref{lem:three-operator-theorems-equivalence} below, all three theorems are equivalent. Furthermore, after we have developed the machinery of adjoints in Section~\ref{sec:adjoint-operators}, we will be able to prove the bounded inverse theorem for Hilbert spaces. Since most of our applications will use Hilbert spaces, we will limit ourselves to the proof for this case and ask the reader to refer to the literature for proofs of the more general cases.

\begin{Lemma}
\label{lem:three-operator-theorems-equivalence}
The open mapping theorem, the bounded inverse theorem, and the closed graph theorem are equivalent.
\end{Lemma}

The proof we give here is good only if $X$ is a Hilbert space. However, the first three parts of the proof apply directly to Banach spaces, and the final part can be generalized by using equivalence classes. (This is a standard technique for getting around the lack of a projection theorem in Banach spaces, but one we won't go into in this text.)

\begingroup
\postdisplaypenalty=10000
\renewcommand{\endProof}{\par\nobreak\hfill$\blacksquare$\par\vspace{0.1cm}}
\begin{Proof}
\textbf{Open mapping theorem $\Rightarrow$ bounded inverse theorem.} It is immediately clear from the hypotheses of the bounded inverse theorem that a linear inverse operator $A^{-1}$ with domain $\mathcal{R}(A)$ exists. The nontrivial assertion is that $A^{-1}$ is bounded. However, this follows from the open mapping theorem, the equivalence of boundedness and continuity for linear operators, and the topological version of the definition of continuity: that an operator $T$ is continuous if and only if the inverse image of open sets in $\mathcal{R}(T)$ is open in $\mathcal{D}(T)$. (The inverse image of an open set in $\mathcal{R}(A^{-1})$ $(=X=\mathcal{D}(A))$ under $A^{-1}$ is the same as the image of the set under $A$.)

\textbf{Bounded inverse theorem $\Rightarrow$ closed graph theorem.} We first observe that the product space $X\times Y$ is a Banach space with norm
\begin{equation}\label{eq:banach-product-sum-norm}
\|(x,y)\|=\|x\|+\|y\|.
\end{equation}
Our hypothesis is that $\Gamma(A)$ is a closed subspace in $X\times Y$ and $\mathcal{D}(A)$ is a closed subspace in $X$. Thus, $\Gamma(A)$ and $\mathcal{D}(A)$ are Banach spaces. We now define a projection map
\begin{equation}\label{eq:graph-domain-projection-map}
P:\Gamma(A)\to\mathcal{D}(A)
\end{equation}
by
\begin{equation}\label{eq:graph-domain-projection-action}
P(x,Ax):=x.
\end{equation}
Note that $P$ is linear and bijective. In fact, its inverse
\begin{equation}\label{eq:graph-domain-projection-inverse-map}
P^{-1}:\mathcal{D}(A)\to\Gamma(A)
\end{equation}
is defined by
\begin{equation}\label{eq:graph-domain-projection-inverse-action}
P^{-1}x:=(x,Ax).
\end{equation}
The mapping $P$ is also bounded since
\begin{equation}\label{eq:graph-domain-projection-boundedness}
\|P(x,Ax)\|=\|x\|\leq\|x\|+\|Ax\|=\|(x,Ax)\|.
\end{equation}
Thus, by the bounded inverse theorem (Theorem~\ref{thm:bounded-inverse}) there is a constant $C$ such that
\begin{equation}\label{eq:bounded-graph-projection-inverse-estimate}
\|(x,Ax)\|=\|P^{-1}x\|\leq C\|x\|.
\end{equation}
But this implies $A$ is bounded since
\begin{equation}\label{eq:closed-graph-operator-boundedness-estimate}
\|Ax\|\leq\|(x,Ax)\|\leq C\|x\|
\end{equation}
for every $x\in\mathcal{D}(A)$.

\textbf{Closed graph theorem $\Rightarrow$ bounded inverse theorem.} This part is left as an exercise. (Problem~\ref{ex:closed-graph-implies-bounded-inverse}.)

\textbf{Bounded inverse theorem $\Rightarrow$ open mapping theorem.} We prove this only in the case where $X$ is a Hilbert space. Since $A$ is bounded, $\mathcal{N}(A)$ is closed (cf. Problem~\ref{ex:closed-operator-null-space-closed}). Thus, we can use the projection theorem to decompose $X$ into $X=\mathcal{N}(A)\oplus\mathcal{N}(A)^\perp$. We then let $P:X\to\mathcal{N}(A)^\perp$ be the orthogonal projection operator and define $\widetilde A$ to be the restriction of $A$ to the domain $\mathcal{N}(A)^\perp$. Observe that $A$ can be written as the composition of these two operators; i.e.,\[Ax=\widetilde A(Px)\]for every $x\in X$. The proof now hinges on two facts which we ask the reader to verify.\begin{enumerate}\item The projection map $P$ maps open sets in $X$ to open sets in $\mathcal{N}(A)^\perp$ (Problem~\ref{ex:orthogonal-projection-open-map}).\item The operator $\widetilde A$ is a continuous bijection from $\mathcal{N}(A)^\perp$ to $Y$ (Problem~\ref{ex:surjective-hilbert-operator-null-complement-bijection}).\end{enumerate}Now, an open set in $X$ gets mapped by $P$ to an open set in $\mathcal{N}(A)^\perp$, and by the bounded inverse theorem, this set gets mapped by $\widetilde A$ to an open set in $Y$. (The image of a set under $\widetilde A$ is the inverse image of a set under $\widetilde A^{-1}$.) Hence, the map $A$, which is the composition of the two maps, takes open sets to open sets.
\end{Proof}
\endgroup
